\section{Significance of Uncertainty Quantification}

A natural question to ask is: why is uncertainty quantification important, and how does it help us to make better decisions?


I want to hereby discuss the type of problems that uncertainty quantification can help us to solve, and hence the type of problems we focus on in this thesis.

In Benjamin Recht's \cite{recht2025uncertainty} blog post series on action-impact and forecasting, he showed that the faster the decision frequency, the less of a role uncertainty quantification played. In cases where feedback is reflected in the observations quickly, the feedback control itself mitigates the impact of uncertainty. This explains why certainty equivalent methods are effective in many 



"Better uncertainty quantification is not going to improve emergency surgeries. It’s not going to improve macroeconomic policy either." These statements seem too broad to me. To say that better uncertainty quantification can't improve settings like macroeconomic policy making sounds like you're saying there is no value in considering how predicted distributions across different models compare. But my experience with policy makers (such as with central bankers, who have invited me to their conferences several times based on their interests in expressing uncertainty) is that decision makers often perceive value in attempts to quantify uncertainty, even if they know the assumptions behind any particular quantification can't be verified. I.e., the problems we tend to see with uncertainty quantification in some of these contexts are not that uncertainty quantification isn't useful or we shouldn't be trying to improve on our methods, it's that some people expect the "small world" view of uncertainty that we can quantify within some model to capture all of our uncertainty. So I wouldn't say better uncertainty quantification is necessarily useless in these settings.

